\documentclass[11pt]{article}
\usepackage{fontspec}
\setmainfont[
  Path=fonts/,
  Extension=.ttf,
  UprightFont=*-400,
  BoldFont=*-700
]{NotoSerifSC}
\setsansfont[
  Path=fonts/,
  Extension=.ttf,
  UprightFont=*-400,
  BoldFont=*-700
]{NotoSerifSC}
\XeTeXlinebreaklocale "zh"
\XeTeXlinebreakskip=0pt plus 1pt
\usepackage[a4paper,margin=1.70cm]{geometry}
\usepackage{amsmath,amssymb,mathtools,bm}
\usepackage{booktabs,longtable,array}
\usepackage{enumitem}
\setlength{\parindent}{2em}
\setlength{\parskip}{0.35em}
\setlist{nosep,leftmargin=2em}
\allowdisplaybreaks[3]
\numberwithin{equation}{section}
\pagestyle{plain}

\newcommand{\R}{\mathbb{R}}
\newcommand{\E}{\mathbb{E}}
\newcommand{\Pp}{\mathbb{P}}
\newcommand{\1}{\mathbf{1}}
\newcommand{\T}{\mathsf{T}}
\newcommand{\F}{\mathrm{F}}
\newcommand{\cN}{\mathcal{N}}
\newcommand{\cL}{\mathcal{L}}
\newcommand{\cS}{\mathcal{S}}
\newcommand{\cH}{\mathcal{H}}
\newcommand{\cA}{\mathcal{A}}
\newcommand{\cM}{\mathcal{M}}
\newcommand{\cT}{\mathcal{T}}
\newcommand{\diag}{\operatorname{diag}}
\newcommand{\tr}{\operatorname{tr}}
\newcommand{\Var}{\operatorname{Var}}
\newcommand{\Cov}{\operatorname{Cov}}
\newcommand{\softmax}{\operatorname{softmax}}
\newcommand{\KL}{D_{\mathrm{KL}}}
\newcommand{\sg}{\operatorname{sg}}
\begin{document}
    \begin{center}
      {\LARGE\bfseries 6-2405.04434v5-Integrated-DeepseekV2}
    \end{center}
    \vspace{0.8em}

    \section{符号与维度}
    \begin{longtable}{>{$}l<{$} >{$}l<{$} p{10cm}}
    \toprule
    \text{符号} & \text{维度} & \text{含义}\\
    \midrule
    X & \R^{T\times d} & 长度为 $T$ 的隐藏状态序列。\\
Q,K,V & \R^{T\times d_h} & 查询、键和值。\\
A & \R^{T\times T} & 行归一化注意力矩阵。\\
M & \R^{T\times T} & 因果或结构掩码。\\
H,d_h & \mathbb N & 注意力头数与单头维度，$Hd_h=d$。\\
c_t^{KV} & \R^{d_c} & MLA 压缩 KV 潜变量。\\
        E_i,g_i & -- & 第 $i$ 个专家与路由分数。\\
    \bottomrule
    \end{longtable}

    所有向量默认是列向量；未特别说明时，期望同时对数据、噪声和采样时刻取值。下文只展开论文中承担方法逻辑的公式，并把所需假设写在推导发生的位置。

\section{缩放点积注意力的逐项推导}

设单头输入为 $X\in\R^{n\times d}$，投影矩阵为
$W_Q,W_K\in\R^{d\times d_k}$、$W_V\in\R^{d\times d_v}$。于是
\begin{equation}
 Q=XW_Q,\qquad K=XW_K,\qquad V=XW_V,
 \qquad S=\frac{QK^\T}{\sqrt{d_k}},
 \qquad P=\operatorname{softmax}_{\rm row}(S),
 \qquad O=PV.
\end{equation}
维度依次为 $Q,K\in\R^{n\times d_k}$、$S,P\in\R^{n\times n}$、
$V,O\in\R^{n\times d_v}$。

\subsection{为什么除以 $\sqrt{d_k}$}

令某一对 query、key 的坐标独立、均值为零、方差为 $\sigma^2$，则
\begin{align}
 q^\T k&=\sum_{a=1}^{d_k}q_ak_a,\\
 \E[q^\T k]&=\sum_a\E[q_a]\E[k_a]=0,\\
 \Var(q^\T k)
 &=\sum_a\Var(q_ak_a)
 =\sum_a\E[q_a^2]\E[k_a^2]
 =d_k\sigma^4.
\end{align}
因此未缩放 logit 的标准差为 $\sqrt{d_k}\sigma^2$；缩放后
\begin{equation}
 \Var\!\left(\frac{q^\T k}{\sqrt{d_k}}\right)=\sigma^4.
\end{equation}
这不是让注意力必然均匀，而是避免初始化时维度增大直接把 softmax 推入饱和区。

\subsection{softmax Jacobian 的完整形式}

对一行 $p_i=e^{s_i}/Z$，$Z=\sum_j e^{s_j}$。逐坐标求导得
\begin{align}
 \frac{\partial p_i}{\partial s_j}
 &=\frac{\delta_{ij}e^{s_i}Z-e^{s_i}e^{s_j}}{Z^2}\\
 &=p_i(\delta_{ij}-p_j).
\end{align}
写成矩阵即
\begin{equation}
 J_{\rm sm}(p)=\diag(p)-pp^\T.
\end{equation}
它是半正定矩阵，因为任意 $z$ 满足
\begin{align}
 z^\T J_{\rm sm}(p)z
 &=\sum_i p_i z_i^2-\left(\sum_i p_i z_i\right)^2\\
 &=\Var_{i\sim p}(z_i)\ge0.
\end{align}
又因 $J_{\rm sm}(p)\1=0$，给所有 logit 加同一常数不改变 softmax。

设损失对输出一行的梯度为 $g_o\in\R^{d_v}$。由 $o=\sum_jp_jv_j$，
\begin{align}
 \frac{\partial\cL}{\partial p_j}&=g_o^\T v_j,\\
 \frac{\partial\cL}{\partial s_j}
 &=\sum_i\frac{\partial\cL}{\partial p_i}
       \frac{\partial p_i}{\partial s_j}\\
 &=p_j g_o^\T(v_j-o).
\end{align}
所以一个 key 的梯度由它相对当前加权均值 $o$ 的偏差决定；这给出了注意力各位置梯度耦合的精确来源。

\subsection{对 $Q,K,V$ 的矩阵梯度}

记 $G_O=\partial\cL/\partial O$。矩阵微分
\begin{equation}
 dO=dP\,V+P\,dV
\end{equation}
与 Frobenius 内积 $d\cL=\langle G_O,dO\rangle_\F$ 给出
\begin{equation}
 G_V=P^\T G_O,
 \qquad
 G_P=G_O V^\T.
\end{equation}
对每行应用 softmax Jacobian，令
\begin{equation}
 G_S=P\odot\left(G_P-\bigl((G_P\odot P)\1\bigr)\1^\T\right),
\end{equation}
其中括号中的列向量在行内广播。再由
\begin{equation}
 dS=\frac{dQK^\T+QdK^\T}{\sqrt{d_k}}
\end{equation}
得到
\begin{equation}
 \boxed{
 G_Q=\frac{G_SK}{\sqrt{d_k}},\qquad
 G_K=\frac{G_S^{\T}Q}{\sqrt{d_k}},\qquad
 G_V=P^{\T}G_O.}
\end{equation}
最后通过 $Q=XW_Q$ 等式继续链式求导：
\begin{align}
 G_{W_Q}&=X^{\T}G_Q,&G_X^{(Q)}&=G_QW_Q^\T,\\
 G_{W_K}&=X^{\T}G_K,&G_X^{(K)}&=G_KW_K^\T,\\
 G_{W_V}&=X^{\T}G_V,&G_X^{(V)}&=G_VW_V^\T.
\end{align}
同一输入同时产生三条梯度路径，故
$G_X=G_X^{(Q)}+G_X^{(K)}+G_X^{(V)}$。

\section{掩码、稀疏化与复杂度}

令二值可见性矩阵 $M\in\{0,1\}^{n\times n}$。严格掩码写成
\begin{equation}
 \widetilde S_{ij}=
 \begin{cases}
 S_{ij},&M_{ij}=1,\\
 -\infty,&M_{ij}=0,
 \end{cases}
 \qquad
 \widetilde P_{ij}=\frac{M_{ij}e^{S_{ij}}}
 {\sum_\ell M_{i\ell}e^{S_{i\ell}}}.
\end{equation}
因此被掩码项的概率与梯度都严格为零。若第 $i$ 行只保留 $s_i$ 个 key，
注意力乘法量为
\begin{equation}
 \Theta\!\left(d_k\sum_i s_i+d_v\sum_i s_i\right).
\end{equation}
稠密情形 $s_i=n$ 为 $\Theta(n^2(d_k+d_v))$；若 $s_i\le s\ll n$，
则为 $\Theta(ns(d_k+d_v))$。但只有稀疏索引和 kernel 真正跳过零块时，理论乘法量才会转化为实际加速。

\subsection{稀疏截断的输出误差}

对一行概率 $p$，保留集合 $S$，被删除概率质量
$\varepsilon=\sum_{j\notin S}p_j$。重新归一化后
\begin{equation}
 \widehat p_j=
 \begin{cases}
 p_j/(1-\varepsilon),&j\in S,\\
 0,&j\notin S.
 \end{cases}
\end{equation}
直接计算其 $\ell_1$ 距离：
\begin{align}
 \|p-\widehat p\|_1
 &=\sum_{j\in S}p_j\frac{\varepsilon}{1-\varepsilon}
   +\sum_{j\notin S}p_j\\
 &=\varepsilon+\varepsilon=2\varepsilon.
\end{align}
若 $\|v_j\|_2\le V_{\max}$，则输出误差满足
\begin{equation}
 \|o-\widehat o\|_2
 =\left\|\sum_j(p_j-\widehat p_j)v_j\right\|_2
 \le 2\varepsilon V_{\max}.
\end{equation}
这说明 top-$k$ 的关键不是 $k$ 本身，而是被丢弃的注意力质量。

\subsection{KV 缓存的精确规模}

自回归解码第 $t$ 步只新增一个 key/value。若有 $L$ 层、$h_{kv}$ 个 KV 头、
每头维度 $d_h$、每元素 $b$ 字节，则缓存量为
\begin{equation}
 \operatorname{Mem}_{KV}=2Lth_{kv}d_hb.
\end{equation}
其中因子 $2$ 来自 key 与 value。多查询或分组查询通过令
$h_{kv}<h_q$ 降低缓存；它改变的是存储与带宽，不改变 query 头的数量。

\section{残差路径与归一化的稳定性}

对 pre-norm 块
\begin{equation}
 x_{\ell+1}=x_\ell+F_\ell(\operatorname{LN}(x_\ell)),
\end{equation}
Jacobian 为
\begin{equation}
 \frac{\partial x_{\ell+1}}{\partial x_\ell}
 =I+J_{F_\ell}J_{\rm LN}.
\end{equation}
跨 $L$ 层梯度是这些矩阵的乘积。与没有残差时只乘
$J_{F_\ell}J_{\rm LN}$ 相比，恒等项提供直接通路。若
$\|J_{F_\ell}J_{\rm LN}\|_2\le\eta_\ell$，则
\begin{equation}
 \left\|\prod_{\ell=1}^{L}(I+A_\ell)\right\|_2
 \le\prod_{\ell=1}^{L}(1+\eta_\ell)
 \le\exp\!\left(\sum_{\ell=1}^{L}\eta_\ell\right).
\end{equation}
这给出稳定性的充分条件；它不是“残差必然防止爆炸”的无条件结论。

对向量 $x\in\R^d$，忽略可学习仿射项时
\begin{equation}
 \mu=\frac1d\1^{\T}x,\qquad
 \sigma^2=\frac1d\|x-\mu\1\|_2^2,\qquad
 y=\frac{x-\mu\1}{\sqrt{\sigma^2+\epsilon}}.
\end{equation}
令 $P_c=I-\frac1d\1\1^\T$、$u=P_cx$、$s=\sqrt{d^{-1}u^{\T}u+\epsilon}$，
则逐步微分得
\begin{align}
 du&=P_c\,dx,\\
 ds&=\frac{u^{\T}du}{d s},\\
 dy&=\frac1sP_c\,dx-\frac{u\,u^{\T}P_c\,dx}{d s^3},
\end{align}
故
\begin{equation}
 J_{\rm LN}=\frac1sP_c-\frac{uu^\T}{d s^3}.
\end{equation}
该式显示归一化会消除全常数方向，并对样本方差方向施加额外投影。

\section{多头分解与合并}

第 $a$ 个头输出 $O_a=P_aV_a\in\R^{n\times d_h}$，拼接后
\begin{equation}
 Y=[O_1\|\cdots\|O_h]W_O,
 \qquad W_O\in\R^{hd_h\times d}.
\end{equation}
因此多头不是把概率矩阵简单平均，而是在不同投影子空间中分别形成概率矩阵，
再由 $W_O$ 线性混合。参数量（忽略偏置）为
\begin{equation}
 d(hd_k)+d(hd_k)+d(hd_v)+(hd_v)d.
\end{equation}
常用 $hd_k=hd_v=d$ 时等于 $4d^2$。注意力矩阵的主内存项仍为
$h n^2$；Flash 类实现通过分块在线维护 softmax 的行最大值与归一化和，避免物化完整矩阵，
但其数学输出仍与上述 softmax 等价。

\section{自回归似然与 token 梯度}

对序列 $x_{1:T}$，因果语言模型按概率链式法则分解
\begin{equation}
 p_\theta(x_{1:T})=\prod_{t=1}^{T}p_\theta(x_t\mid x_{<t}).
\end{equation}
负对数似然
\begin{equation}
 \cL_{\rm NLL}=-\sum_{t=1}^{T}\log p_\theta(x_t\mid x_{<t}).
\end{equation}
若第 $t$ 步 logits 为 $z_t\in\R^{V}$，概率
$p_{tj}=e^{z_{tj}}/\sum_ke^{z_{tk}}$，真实 token 索引 $y_t$，则
\begin{equation}
 \ell_t=-z_{t,y_t}+\log\sum_ke^{z_{tk}},
\end{equation}
逐 logit 梯度
\begin{equation}
 \frac{\partial\ell_t}{\partial z_{tj}}
 =p_{tj}-\1\{j=y_t\}.
\end{equation}
Hessian 正是 softmax Jacobian
\begin{equation}
 \nabla_z^2\ell_t=\diag(p_t)-p_tp_t^T\succeq0.
\end{equation}
所以交叉熵关于 logits 凸，但关于深网参数并不凸。

对非 padding token 集合 $\mathcal I$，平均损失和困惑度
\begin{equation}
 \bar L=\frac1{|\mathcal I|}\sum_{t\in\mathcal I}\ell_t,
 \qquad \operatorname{PPL}=e^{\bar L}.
\end{equation}
因此 PPL 是真实 token 概率倒数的几何平均：
\begin{equation}
 \operatorname{PPL}=\left(
 \prod_{t\in\mathcal I}\frac1{p_\theta(x_t\mid x_{<t})}
 \right)^{1/|\mathcal I|}.
\end{equation}
不同 tokenizer 的 token 粒度不同，PPL 不能不经换算直接横向比较。

\section{旋转位置编码的相对位置恒等式}

对二维坐标对，定义旋转
\begin{equation}
 R(m\omega)=
 \begin{bmatrix}
 \cos(m\omega)&-\sin(m\omega)\\
 \sin(m\omega)&\cos(m\omega)
 \end{bmatrix}.
\end{equation}
位置 $m,n$ 的 query、key 分别为
$q_m=R(m\omega)q$、$k_n=R(n\omega)k$。内积
\begin{align}
 q_m^Tk_n
 &=q^TR(m\omega)^TR(n\omega)k\\
 &=q^TR((n-m)\omega)k.
\end{align}
第二步用到 $R(a)^T=R(-a)$ 和 $R(a)R(b)=R(a+b)$。
所以绝对位置旋转后的内积只依赖相对距离 $n-m$。

在 $d_h$ 维中对相邻坐标对使用频率
\begin{equation}
 \omega_j=\theta^{-2j/d_h},\qquad j=0,\ldots,d_h/2-1.
\end{equation}
低 $j$ 频率高、分辨近距离，高 $j$ 频率低、覆盖远距离。若推理长度远超训练长度，
相位 $m\omega_j$ 的外推和周期混叠会改变注意力几何；位置插值本质上是重新缩放这些相位。

\section{RMSNorm 的 Jacobian}

忽略可学习增益，RMSNorm
\begin{equation}
 y=\frac{x}{s},\qquad
 s=\sqrt{\frac1d x^Tx+\epsilon}.
\end{equation}
微分
\begin{equation}
 ds=\frac{x^Tdx}{d s},
\end{equation}
因此
\begin{align}
 dy&=\frac1s dx-\frac{x}{s^2}ds\\
 &=\left(\frac1sI-\frac{xx^T}{d s^3}\right)dx.
\end{align}
故
\begin{equation}
 J_{\rm RMS}=\frac1sI-\frac{xx^T}{d s^3}.
\end{equation}
与 LayerNorm 相比，它没有去均值投影
$I-d^{-1}\1\1^T$，因此不消除常数方向。加上逐坐标增益 $g$ 后，
Jacobain 左乘 $\diag(g)$。

谱上，对任意 $v\perp x$，$Jv=v/s$；沿 $x$ 方向
\begin{equation}
 Jx=\left(\frac1s-\frac{\|x\|^2}{d s^3}\right)x
 =\frac{\epsilon}{s^3}x.
\end{equation}
当 $\epsilon$ 很小，径向尺度方向被强烈压制，而切向方向按 $1/s$ 缩放。

\section{SwiGLU 前馈层的逐项梯度}

SwiGLU 可写
\begin{equation}
 a=xW_a,\qquad b=xW_b,\qquad
 h=\operatorname{swish}(a)\odot b,\qquad
 y=hW_o,
\end{equation}
其中
\begin{equation}
 \operatorname{swish}(a)=a\sigma(a).
\end{equation}
其导数
\begin{align}
 \frac d{da}[a\sigma(a)]
 &=\sigma(a)+a\sigma(a)(1-\sigma(a)).
\end{align}
令 $g_h=\partial\cL/\partial h$，则
\begin{align}
 g_a&=g_h\odot b\odot
 [\sigma(a)+a\sigma(a)(1-\sigma(a))],\\
 g_b&=g_h\odot\operatorname{swish}(a),\\
 \nabla_{W_a}\cL&=x^Tg_a,\qquad
 \nabla_{W_b}\cL=x^Tg_b,\\
 \nabla_x\cL&=g_aW_a^T+g_bW_b^T.
\end{align}
门分支 $b$ 同时调制激活和 $a$ 分支梯度，故两个投影不是可独立解释的普通 MLP。

\section{MoE 路由与负载}

对 token $x_i$，router logits $r_i=W_rx_i$，概率
\begin{equation}
 p_{ie}=\frac{e^{r_{ie}}}{\sum_{j=1}^{E}e^{r_{ij}}}.
\end{equation}
若 top-$k$ 专家集合 $S_i$，归一化门
\begin{equation}
 \widetilde p_{ie}=\frac{p_{ie}\1\{e\in S_i\}}
 {\sum_{j\in S_i}p_{ij}},
 \qquad
 y_i=\sum_{e\in S_i}\widetilde p_{ie}F_e(x_i).
\end{equation}
固定集合 $S_i$ 内可正常求导；集合边界的 arg-top-$k$ 不连续，几乎处处把未选专家梯度置零。

令软重要性
\begin{equation}
 P_e=\frac1N\sum_ip_{ie},
\end{equation}
硬负载
\begin{equation}
 F_e=\frac1N\sum_i\1\{e\in S_i\}/k.
\end{equation}
常见平衡项
\begin{equation}
 L_{\rm bal}=E\sum_eP_eF_e
\end{equation}
在均匀 $P_e=F_e=1/E$ 时等于 $1$。把硬 $F_e$ 停止梯度后，
router 梯度只通过 $P_e$，会降低当前过载专家的概率；它是训练启发式，不等同于严格容量约束。

若每专家容量为
\begin{equation}
 C_e=\left\lceil\frac{Nk}{E}\,c_{\rm cap}\right\rceil,
\end{equation}
超过容量的 token 必须丢弃或转送。即使平均负载均匀，局部批次方差仍可能溢出。

\section{低秩 KV 潜变量的吸收}

设 key/value 由共享潜变量
\begin{equation}
 c_t=W_cx_t\in\R^{d_c},\qquad
 k_t=W_kc_t,\qquad v_t=W_vc_t.
\end{equation}
注意力 logit
\begin{equation}
 q_s^Tk_t=q_s^TW_kc_t=(W_k^Tq_s)^Tc_t.
\end{equation}
因此解码时可把 $W_k$ 吸收到 query 投影，缓存 $c_t$ 而非完整 $k_t$。
若 value 聚合
\begin{equation}
 o_s=\sum_tp_{st}W_vc_t=W_v\left(\sum_tp_{st}c_t\right),
\end{equation}
$W_v$ 可放到加权和之后。缓存从每 token 的 $d_k+d_v$ 降为 $d_c$，
前提是位置编码等操作不破坏这些线性交换关系；对只作用在部分 key 维度的 RoPE，需把旋转部分另行缓存或重构。

\section{多 token 预测的联合似然}

若同一隐藏状态 $h_t$ 用 $K$ 个头预测 $x_{t+1},\ldots,x_{t+K}$，
\begin{equation}
 \cL_t=\sum_{k=1}^{K}\lambda_k
 [-\log p_{\theta,k}(x_{t+k}\mid h_t)].
\end{equation}
共享骨干梯度
\begin{equation}
 \nabla_{h_t}\cL_t=\sum_{k=1}^{K}\lambda_k g_{t,k}.
\end{equation}
其平方范数精确展开
\begin{equation}
 \left\|\sum_k\lambda_kg_{t,k}\right\|^2
 =\sum_k\lambda_k^2\|g_{t,k}\|^2
 +2\sum_{j<k}\lambda_j\lambda_k g_{t,j}^Tg_{t,k}.
\end{equation}
交叉内积为负时存在任务冲突；辅助头退火相当于训练后期逐渐去掉这些交叉项，
使最终目标回到下一 token 似然。

\section{MLA 的低秩 KV 潜变量}
标准多头注意力缓存每个 token 的 $K,V\in\R^{H\times d_h}$。MLA 先压缩
\begin{equation}
  c_t^{KV}=W^{DKV}h_t\in\R^{d_c},
\end{equation}
再按头上投影
\begin{equation}
  k_{t,h}^C=W_h^{UK}c_t^{KV},\qquad
  v_{t,h}=W_h^{UV}c_t^{KV}.
\end{equation}
缓存从每 token 约 $2Hd_h$ 个数降为 $d_c$ 加少量位置键。该分解把 KV 映射秩限制在 $d_c$。

\section{吸收矩阵与解码}
查询内容部分 $q_h^C=W_h^{UQ}c_t^Q$ 与键的点积
\begin{equation}
  (q_h^C)^Tk_{j,h}^C
  =(c_t^Q)^T(W_h^{UQ})^TW_h^{UK}c_j^{KV}.
\end{equation}
可预先吸收中间矩阵，避免显式恢复完整键。位置部分单独使用 RoPE，从而不破坏可吸收的低秩内容通道。

\section{细粒度 MoE}
若 $N_r$ 个共享专家始终激活、从 $N_s$ 个路由专家选 top-$K$，输出
\begin{equation}
  y=\sum_{i\in\cS_{\rm shared}}E_i(x)
  +\sum_{i\in\operatorname{TopK}(g(x))}p_iE_i(x).
\end{equation}
每 token 计算仅与 $N_r+K$ 成正比，而参数容量随 $N_r+N_s$ 增长。负载均衡损失用于防止路由塌缩，但会与任务最优路由产生权衡。

\end{document}
